\chapter{Inteligencia de Amenazas (Threat Intelligence)}
\label{ch:threat_intelligence}

En este capítulo se estructura el modelo de vigilancia activa de amenazas para anticipar ataques dirigidos contra el entorno universitario.

\section{Fuentes y Vectores de Vigilancia}
El programa de Inteligencia de Amenazas de la UNAP se enfoca en el monitoreo continuo de los siguientes riesgos críticos:
\begin{itemize}
    \item \textbf{Phishing de Credenciales Institucionales:} Campañas dirigidas a capturar cuentas de correo \texttt{@unap.edu.pe} y accesos a matrícula.
    \item \textbf{Dominios Tipográficos (Typosquatting):} Detección temprana de registros maliciosos que imitan los portales oficiales (ej. dominios apócrifos como \texttt{unap-matricula.*}, \texttt{unap-portal.*}).
    \item \textbf{Vulnerabilidades y Exploits Conocidos:} Alertas sobre fallos que afectan al stack tecnológico observado, en especial exploits dirigidos a \texttt{Nginx 1.16.1} y componentes web legados.
    \item \textbf{Fuentes de Información Especializadas:} Monitoreo de feeds del CSIRT Nacional (Perú), catálogo CISA KEV (Known Exploited Vulnerabilities) y comunidades globales de inteligencia (AlienVault OTX).
\end{itemize}

\section{Procedimiento de Gestión de Inteligencia}
El flujo operativo para transformar datos de amenazas en contramedidas defensivas sigue un ciclo de 6 etapas:

\begin{center}
\small
\textbf{Recopilación} $\longrightarrow$ \textbf{Correlación Logs} $\longrightarrow$ \textbf{Clasificación Criticidad} \\
$\longrightarrow$ \textbf{Generación IoCs} $\longrightarrow$ \textbf{Actualización WAF} $\longrightarrow$ \textbf{Comunicación Interna}
\end{center}

\begin{enumerate}
    \item \textbf{Recopilación:} Captura de feeds de IoCs (direcciones IP maliciosas, hashes de malware, dominios fraudulentos).
    \item \textbf{Correlación de Logs:} Búsqueda retrospectiva en la plataforma SIEM para verificar si los IoCs tuvieron contacto con la red interna.
    \item \textbf{Clasificación de Criticidad:} Evaluación del impacto potencial de la amenaza reportada sobre los activos de la UNAP.
    \item \textbf{Generación de IoCs:} Creación de firmas y listas negras institucionales.
    \item \textbf{Actualización en Firewall/WAF:} Ingesta automática de bloqueos en el firewall perimetral y reglas del WAF.
    \item \textbf{Comunicación Interna:} Emisión de alertas tempranas y boletines preventivos dirigidos a los usuarios.
\end{enumerate}

\section{Producto de la Sección: Boletín de Inteligencia de Amenazas}

\begin{cmdbox}{Boletín de Inteligencia de Amenazas: Alerta Spear-Phishing y Suplantación}
[ALERTA DE SEGURIDAD DIGITAL - CSIRT-UNAP / OTI]
Codigo de Alerta: ALERTA-TI-2026-004
Nivel de Severidad: ALTO (CVSS: 7.8)
Fecha de Emision: 17 de Octubre de 2026

AMENAZA DETECTADA:
Campana de Spear-Phishing dirigida a docentes y alumnos mediante dominios de
typosquatting que emulan el sistema oficial de matricula (ej. inap-matricula.com).

INDICADORES DE COMPROMISO (IoCs):
* Dominios detectados: unap-matricula.online, portal-unap.org
* IPs de origen C2: 198.51.100.45, 203.0.113.88
* Asunto de correos fraudulentos: 'URGENTE: Actualizacion de Datos Matricula 2026-I'

ACCIONES PREVENTIVAS EJECUTADAS:
1. Bloqueo de dominios e IPs en Firewall de borde y servidores DNS internos.
2. Ingesta de reglas de deteccion en Wazuh SIEM para correlacionar accesos.
3. Notificacion a los administradores de correo para aislamiento de mensajes.
\end{cmdbox}